% Part: intuitionistic-logic
% Chapter: tableaux
% Section: rules

\documentclass[../../../include/open-logic-section]{subfiles}

\begin{document}

\olfileid{int}{tab}{rul}

\olsection{Aturan kanggo Logika Intuisionistik}

Aturan kanggo konektif $\land$ lan $\lor$ padha karo kanggo
!!{tableau}s proposisional mawa tandha biasa, mung ditambahi prefiks.
Ing saben kasus, aturan kang ditrapake marang !!{formula} mawa tandha
$\sFmla{S}{!A}[\sigma]$ ngasilake !!{formula}s anyar kang uga
mawa prefiks~$\sigma$. Iki cetha sacara intuitif: upamane yen
$!A \land !B$ bener ing (donya kang dijenengi dening)~$\sigma$,
mula $!A$ lan $!B$ bener ing~$\sigma$ (aturan iki ngenalake
formula ing prefiks iku, tanpa nyingkirake kabeneran ing donya liyane).
Aturan kanggo $\land$ lan $\lor$ diklumpukake ing \olref{tab:prop-rules}.

\begin{table}
  \[\def\arraystretch{3}\begin{array}{|c|c|}
    \hline
    \AxiomC{\sFmla{\True}{!A \land !B}[\sigma]}
    \RightLabel{\TRule{\True}{\land}}
    \UnaryInfC{\sFmla{\True}{!A}[\sigma]}
    \noLine
    \UnaryInfC{\sFmla{\True}{!B}[\sigma]}
    \DisplayProof
    &
    \AxiomC{\sFmla{\False}{!A \land !B}[\sigma]}
    \RightLabel{\TRule{\False}{\land}}
    \UnaryInfC{$\sFmla{\False}{!A}[\sigma] \quad \mid \quad
      \sFmla{\False}{!B}[\sigma]$}
    \DisplayProof
    \\[2ex]
    \hline
    \AxiomC{\sFmla{\True}{!A \lor !B}[\sigma]}
    \RightLabel{\TRule{\True}{\lor}}
    \UnaryInfC{$\sFmla{\True}{!A}[\sigma] \quad \mid \quad
      \sFmla{\True}{!B}[\sigma]$}
    \DisplayProof
    &
    \AxiomC{\sFmla{\False}{!A \lor !B}[\sigma]}
    \RightLabel{\TRule{\False}{\lor}}
    \UnaryInfC{\sFmla{\False}{!A}[\sigma]}
    \noLine
    \UnaryInfC{\sFmla{\False}{!B}[\sigma]}
    \DisplayProof
    \\[2ex]
    \hline
  \end{array}\]
  \caption{Aturan !!{tableau} mawa prefiks kanggo $\land$ lan $\lor$}
  \ollabel{tab:prop-rules}
\end{table}

Syarat penutupan padha karo kanggo !!{tableau}s biasa, nanging ora
mung !!{formula}s kang kudu cocog, prefikse uga kudu cocog.
Sejatine, syarate bisa luwih longgar: ing model intuisionistik,
!!{formula}s kang wis bener tetep bener, mula ora mungkin $!A$
bener ing~$\sigma$ nanging salah ing prefiks aksesibel
~$\sigma.{*}$. Dadi cabang katutup yen ngemot loro-lorone
\[
\sFmla{\True}{!A}[\sigma] \quad\text{lan}\quad \sFmla{\False}{!A}[\sigma.{*}]
\]
kanggo sawenehing prefiks $\sigma$ lan !!{formula}~$!A$.
Yen tandhane dibalik, yaiku yen ngemot
\[
\sFmla{\False}{!A}[\sigma] \quad\text{lan}\quad \sFmla{\True}{!A}[\sigma.{*}]
\]
cabang katutup mung yen $*$ iku urutan kosong.

Kajaba iku, cabang katutup yen ngemot~$\sFmla{\True}{\bot}[\sigma]$.

Aturan kanggo nyetel asumsi uga padha karo kanggo !!{tableau}s
biasa, nanging kanggo asumsi kita tansah nggunakake prefiks~$1$.
(Ora dadi prakara prefiks endi kang digunakake, anggere padha
kanggo kabeh asumsi.) Dadi, upamane, kita nyebut
\[
!B_1, \dots, !B_n \Proves !A
\]
yen lan mung yen ana tabelau katutup kanggo asumsi
\[
\sFmla{\True}{!B_1}[1], \dots, \sFmla{\True}{!B_n}[1],
\sFmla{\False}{!A}[1].
\]

Kanggo kondisional~$\lif$, aturane beda karo kasus klasik lan modal.
Aturan $\TRule{\True}{\lif}$ ngluwihi cabang kang ngemot
$\sFmla{\True}{!A \lif !B}[\sigma]$ nganggo
$\sFmla{\False}{!A}[\sigma.{*}]$ lan $\sFmla{\True}{!B}[\sigma.{*}]$
ing rong cabang kang beda. Aturan iki mung bisa ditrapake kanggo
prefiks~$\sigma.{*}$ kang \emph{wis} ana ing cabang kang dadi
papan penerapane. Prefiks kaya mangkono diarani ``wis digunakake''
(ing cabang iku). (Amarga $\sigma.{*}$ ngemot $\sigma$ dhewe,
aturan iki tansah bisa ditrapake kanthi nambahake formula mawa
tandha lan prefiks $\sFmla{\False}{!A}[\sigma]$ lan
$\sFmla{\True}{!B}[\sigma]$ ing cabang kang kapisah.)

Aturan $\TRule{\False}{\lif}$ ngluwihi cabang kang ngemot
$\sFmla{\False}{!A \lif !B}[\sigma]$ nganggo loro-lorone
$\sFmla{\True}{!A}[\sigma.n]$ lan $\sFmla{\False}{!B}[\sigma.n]$
ing cabang kang padha, kanthi $\sigma.n$ prefiks anyar ing cabang iku.

Aturan kanggo $\lnot$ ditetepake kanthi cara kang padha
(nggunakake definisi $\lnot !A$ minangka $!A \lif \lfalse$).

Aturan iku diwenehake ing \olref{tab:rules-lif-lnot}.

\begin{table}
  \[\def\arraystretch{3}\begin{array}{|c|c|}
    \hline
    \AxiomC{\sFmla{\True}{\lnot !A}[\sigma]}
    \RightLabel{\TRule{\True}{\lnot}}
    \UnaryInfC{$\sFmla{\False}{!A}[\sigma.{*}]$}
    \DisplayProof
    &
    \AxiomC{\sFmla{\False}{\lnot !A}[\sigma]}
    \RightLabel{\TRule{\False}{\lnot}}
    \UnaryInfC{\sFmla{\True}{!A}[\sigma.n]}
    \DisplayProof
    \\[1ex]
    \text{$\sigma.{*}$ wis digunakake} & \text{$\sigma.n$ anyar}\\
    \hline
    \AxiomC{\sFmla{\True}{!A \lif !B}[\sigma]}
    \RightLabel{\TRule{\True}{\lif}}
    \UnaryInfC{$\sFmla{\False}{!A}[\sigma.{*}] \quad \mid \quad
      \sFmla{\True}{!B}[\sigma.{*}]$}
    \DisplayProof
    &
    \AxiomC{\sFmla{\False}{!A \lif !B}[\sigma]}
    \RightLabel{\TRule{\False}{\lif}}
    \UnaryInfC{\sFmla{\True}{!A}[\sigma.n]}
    \noLine
    \UnaryInfC{\sFmla{\False}{!B}[\sigma.n]}
    \DisplayProof
    \\[1ex]
    \text{$\sigma.{*}$ wis digunakake} & \text{$\sigma.n$ anyar}\\
    \hline
  \end{array}\]
  \caption{Aturan !!{tableau} mawa prefiks kanggo $\lnot$ lan $\lif$}
  \ollabel{tab:rules-lif-lnot}
\end{table}

\end{document}
